% 第16批：独立推导、自查；依赖源图的题标明判读前提。
\subsection{自测题4：振动与波动计算、问答}

\sourcetest{4.3.1}{96}{104}
\problemtext 质点位移 $x=0.24\cos(\pi t/2+\pi/3)$（SI），求从 $t=0$ 到首次达到 $x=-0.12$m、且 $v<0$ 的时间。
\workings 角频率 $\omega=\pi/2$；末态满足 $\cos\theta=-1/2$ 且 $\sin\theta>0$，故最早取 $\theta=2\pi/3$。初相 $\pi/3$，于是 $\boxed{\Delta t=(2\pi/3-\pi/3)/(\pi/2)=2/3\ {\rm s}}$。
\crosscheck 旋转矢量沿正向转过 $\pi/3$；另一个位移根 $4\pi/3$ 对应 $v>0$。

\sourcetest{4.3.2}{96}{104}
\problemtext 质量2kg的物块放于竖直振动平板上，板周期0.5s、振幅4cm；板初始在上方最大位移处，求物块对板压力与开始离板的临界振幅。
\workings 取竖直向上为正、$z=A\cos\omega t$，$\omega=2\pi/T=4\pi$rad/s。接触时 $N-mg=m\ddot z=-m\omega^2 z$，即 $N=m(g-\omega^2 z)$。若板在最高点，$N_{\min}=m(g-\omega^2A)\simeq2(9.8-16\pi^2\cdot0.04)=\boxed{6.97\ {\rm N}}$；最低点 $N_{\max}=m(g+\omega^2A)\simeq\boxed{32.23\ {\rm N}}$。临界最高点 $N_{\min}=0$，$A_c=g/\omega^2\simeq\boxed{6.21\ {\rm cm}}$。
\crosscheck 原题表述“物体对平板的压力”数值等于平板支持力；从最高点出发处压力为最小值，尚未离板。

\sourcetest{4.3.3}{96}{104}
\problemtext 题4.3.1图给出向左传播的同一简谐波在 $t=0$ 和 $t=2$s 的波形，两图向左平移20m；已知原点处 $t=0$ 过零且初速度为正，$t=2$s 位移为 $A/\sqrt2$，求原点振动式和波动式。
\workings 左行波写为 $y=A\sin(\omega t+kx)$。在原点两时刻相位由0增加 $\pi/4$，故取图示最小相移 $\omega(2\,{\rm s})=\pi/4$，$\omega=\pi/8$rad/s。两图相移20m对应 $k\,20=\pi/4$，$k=\pi/80$rad/m，从而 $u=\omega/k=10$m/s。故 $\boxed{y(0,t)=A\sin(\pi t/8)}$，$\boxed{y(x,t)=A\sin(\pi t/8+\pi x/80)}$。
\crosscheck 波形在2s内向左移20m；若不排除整周期别名 $\omega 2=\pi/4+2n\pi$，单凭两幅快照不能唯一决定频率；本解采用原图绘示的最短相移分支，需保留图示条件标记。

\sourcetest{4.3.4}{96}{104}
\problemtext 一右行波速500m/s，点 $P$ 位于 $x=1$m，$y_P=0.03\cos(500\pi t-\pi/2)$（SI）；求全波波动式及 $t=0$ 波形。
\workings $\omega=500\pi$rad/s，$k=\omega/u=\pi$rad/m。设 $y=0.03\cos(500\pi t-\pi x+\varphi)$；于 $x=1$ 应有 $-\pi+\varphi=-\pi/2$（模 $2\pi$），故 $\varphi=\pi/2$。$\boxed{y(x,t)=0.03\cos(500\pi t-\pi x+\pi/2)}$；$\boxed{y(x,0)=0.03\sin(\pi x)}$。
\crosscheck 验证在 $x=1$m 时恢复题给 P 点振动式；波长2m。

\sourcetest{4.3.5}{96}{104}
\problemtext 右行声波频率 $f_n$、波长 $\lambda_n$，遇到向左匀速移动的刚性平面（速率 $v_0<c$）后反射；求 $\Delta t$ 内入射波峰数、反射波最前端距离以及反射波长。
\workings 入射波与反向运动平面的相遇速率为 $c+v_0$，碰到平面的波峰数 $\boxed{N=f_n(1+v_0/c)\Delta t}$。最初反射波峰左行速率 $c$，平面左移 $v_0\Delta t$，故两者末距 $\boxed{d=(c-v_0)\Delta t}$。反射波峰间距 $\lambda_r=(c-v_0)/(f_n(1+v_0/c))=\boxed{\lambda_n(c-v_0)/(c+v_0)}$。
\crosscheck 此题壁面靠近声源，反射波波长缩短；当 $v_0\to0$ 恢复 $\lambda_r=\lambda_n$。

\sourcetest{4.4.1}{96}{104}
\problemtext 题4.4.1图为向 $+x$ 传播的余弦式简谐波 $t=0$ 的波形；求图中 O、2、3 等标点的初相并说明如何画出 $t=T/4$ 波形。
\workings 图中 $t=0$ 的波形可写作 $y(x,0)=A\sin(kx)$（相邻过零点间距为半波长）；右行波为 $\boxed{y(x,t)=A\sin(kx-\omega t)}$。改写成余弦振动形式 $A\cos(\omega t-kx+\varphi)$，初相在 $O(x=0)$ 为 $\pi/2$，$x=2$（下降过零）为 $-\pi/2$，$x=3$（波谷）为 $-\pi$（均按 $2\pi$ 取模）。$t=T/4$ 波形是原图整体向右平移 $\boxed{\lambda/4}$。
\crosscheck 以题图横轴标注 $x=1,2,3,4$ 理解 $x$ 的单位；原图局部峰谷可通过代入波形验证。
